
This work tested whether coefficient of variation (CV) of linear probe accuracy trajectories on frozen CLIP features could distinguish spurious from core features. The hypothesis was decisively refuted: AUC = 0.0 on Waterbirds, with both feature types showing nearly identical CV values (~0.04).

The root cause is feature saturation in pretrained models. CLIP has already learned both background and bird type concepts. There are no emergence dynamics to measure---probes converge immediately at all regularization strengths.

The key insight is that emergence uniformity is a training-time phenomenon that cannot be recovered from frozen pretrained representations. The analogy: one cannot measure how fast someone learned by examining their final exam score.

If emergence uniformity proves measurable during training with unfrozen representations, it could still enable single-run, annotation-free spurious correlation mitigation. Future work should focus on training-time CV measurement, intermediate layer probing during training, and loss-based emergence signals.

